\section{Introduction}
\subsection{Unique-Games}\label{sec:ug}
The Unique Games Conjecture (UGC in short) 
is a central open problems in Complexity Theory~\cite{Khot02}.
The primary reason for the interest in UGC is that, if true, it
implies a large number of hardness of approximation results 
that are often times tight~\cite{KhotR08,KhotKMO07,Austrin,KNSgroth,Raghavendra08}
(see~\cite{KhotICM,trevisan2012khot}). For example, 
one of the most striking
consequences of UGC is that it 
implies that a class of semi-definite programs
achieves the best possible approximation ratio (among all 
efficient algorithms) for constraint satisfaction problems~\cite{Raghavendra08}.

Research towards a proof of UGC, on the other hand, 
has stalled for the most part up until recently. Unlike problems such 
as $3$-SAT, a randomly sampled instance of Unique-Games (UG) is easy to solve, raising the question of how do hard instances of UG
even look like. To discuss this, we first give 
a formal definition of the Unique-Games problem and the statement of the Unique-Games Conjecture.
\begin{definition}
A instance of Unique-Games $\Psi$ consists of a graph $G = (V,E)$, 
a finite alphabet $\Sigma$ and a collection of constraints, 
$\Phi = \{\Phi_e\}_{e\in E}$, one for each edge in $G$. 
For all $e\in E$, the constraint $\Phi_e$ takes the form 
$\Phi_e = \{(\sigma, \phi_e(\sigma))~|~\sigma\in \Sigma\}$, 
where $\phi_e\colon \Sigma\to\Sigma$ is a $1$-to-$1$ map.

The goal in the Unique-Games problem is to find an assignment 
$A\colon V\to\Sigma$ that satisfies the maximum number of constraints possible, that is, satisfies that $(A(u),A(v))\in \Phi_{e}$ 
for the largest number of edges $e=(u,v)\in E$ as possible. We define
 the value of the instance $\Psi$ by:
\[
{\sf val}(\Psi) = \max_{A\colon V\to\Sigma}\frac{\#\{e~|~\text{A satisfies $e$}\}}{|E|}.
\]
\end{definition}
With this in mind, the Unique-Games Conjecture is the following statement:
\begin{conjecture}
For all $\eps,\delta>0$ there is $k\in\mathbb{N}$ such that 
given a Unique-Games instance $\Psi$ with alphabet size at most $k$,
it is NP-hard to distinguish between:

\textbf{YES case:} ${\sf val}(\Psi)\geq 1-\eps$.

\textbf{NO case:} ${\sf val}(\Psi)\leq \delta$.
\end{conjecture}

It turns out that the topology of the underlying graph $G$ plays a crucial role in the complexity of the UG instance defined over it~\cite{AroraKKSTV08,MMlocalexpander,DBLP:journals/eccc/AroraIMS10}. In particular, it 
turns out that UG over expander graphs is easy:

% It turns out that the topology of the underlying graph $G$ plays a crucial role in the complexity of the Unique-Games instance defined over it. In~\cite{AroraKKSTV08}, it is shown that if 
% $\Psi$ is a Unique-Games instance with value at least $1-\eps$ and the underlying graph $G$ is an expander graph with spectral gap $\lambda$, then 
% one can find in polynomial time an assignment satisfying at least 
% $1-O(\eps/\lambda)$ of the constraints. Thus, solving Unique-Games over expander graphs is computationally easy. As random graphs are expanders it follows that ``typical'' instances of Unique-Games are easy. 

% In~\cite{MMlocalexpander,DBLP:journals/eccc/AroraIMS10}, it is 
% further shown that a weaker and more local combinatorial 
% notion of expansion suffices to make that conclusion, and 
% to state this we define the notion of edge expansion:
\begin{definition}
Given a regular graph $G=(V,E)$ and a set 
of vertices $S\subseteq V$, the edge expansion of $S$ is defined by:
\[
\Phi(S) = \Pr_{u\in S, v\in \Gamma(u)}\left[v\not\in S\right].
\]
\end{definition}
A graph $G$ is a called a $(\gamma,\xi)$-small set-expander if for every $S\subseteq V$ of size at most $\xi |V|$  it holds that $\Phi(S)\geq \gamma$. Informally, we say that $G$ is a small set expander if it is a $(\gamma,\xi)$-small set 
expander for $\xi$ that is a small constant, and we say $G$ is an expander 
if $\xi = 1/2$. The results of~\cite{AroraKKSTV08,MMlocalexpander,DBLP:journals/eccc/AroraIMS10,BBKSS}
 assert that UG instances with completeness close to $1$ over small-set expanders are easy. 
%Recently, this result was improved in~\cite{BBKSS} for a special class of small-set expander graphs which they refer 
% to as ``certifiable small set expanders''.
% \footnote{Roughly speaking, a graph is 
% a certifiable smalls set expander if 
% there is an \sos proof of constant degree 
% that shows it is a small set expander.}.
% More specifically, they showed a 
% non-trivial algorithm for Unique-Games over graphs that are 
% small-set expanders with somewhat
% weak $\gamma$ and $\xi$.
Thus, to have any chance of proving UGC, one must use graphs 
which are outside the scope of expanders and small set expanders.
% A graph $G$ is a called a $(\gamma,\xi)$-small set-expander
% if any $S\subseteq V$ of size at most $\gamma|V|$  it holds that $\Phi(S)\geq 1-\xi$, and informally we say that $G$ is a small set expander if it is a $(\gamma,\xi)$-small set 
% expander for $\gamma,\xi$ that are small constants. For such graphs the results of~\cite{DBLP:journals/eccc/AroraIMS10} 
% assert that given a Unique-Games instance with value at least 
% $1-\eps$ whose underlying graph is $(\gamma,\xi)$-small set-expander, one can efficiently find an assignment satisfying at least $1-O_{\gamma,\xi}(\eps)$ of the constraints. This result requires $\eps$ 
% to be small enough compared to $\gamma$ and $\xi$ for the algorithm to have a meaningful guarantee. This limitation was recently lifted in~\cite{BBKSS} for a special class of small-set expander graphs which they refer 
% to as ``certifiable small set expanders''.
% \footnote{Roughly speaking, a graph is 
% a certifiable smalls set expander if 
% there is an \sos proof of constant degree 
% that shows it is a small set expander.}
% It thus follows 
% that to have any chance of proving UGC, one must use graphs 
% which are outside the scope of expanders and small set expanders. 


Indeed, recent progress towards UGC~\cite{KhotMS17,DBLP:conf/stoc/DinurKKMS18a,DinurKKMS18,KhotMS18}
has utilized graphs which are not small-set expanders. In these 
works it is proved that $2$-to-$1$-Games are NP-hard (which is a 
very similar problem to UG, except that 
each one of the maps $\phi_e$ defining the constraints 
is a $2$-to-$1$ map). Among others, this implies that for all $\eps > 0$, given
a UG instance $\Psi$ over sufficiently large alphabet, 
it is NP-hard to distinguish between the case that 
${\sf val}(\Psi)\geq 1/2$ and the case that 
${\sf val}(\Psi)\leq \eps$. To prove these results, these works
use graphs that are not small set expanders in two different ways:
\begin{enumerate}
    \item Smooth Parallel Repetition: A key step in the reduction     of~\cite{KhotMS17,DBLP:conf/stoc/DinurKKMS18a,DinurKKMS18,KhotMS18} is an application of the Parallel Repetition Theorem~\cite{Raz} to get a hardness result 
    for a sufficiently 
    smooth outer PCP construction. Roughly speaking, this 
    step in the process may be associated with the Johnson graph
    with a large intersection parameter. That is, with the graph 
    $J(n,\ell,t)$ in which the vertices are $\binom{[n]}{\ell}$, and 
    two vertices $A$ and $B$ are adjacent if $|A\cap B|=\ell-t$, 
    and we think of $t$ as much smaller than $\ell$ (say, $t=\sqrt{\ell}$).
    \item Composition with the Grassmann encoding: The Grassmann
    encoding is an encoding of linear functions based on the Grassmann graph 
    ${\sf Grass}(n,\ell)$ over $\mathbb{F}_2$. The Grassmann graph 
    over $\mathbb{F}_2$ is the graph whose vertices are all 
    $\ell$-dimensional subspaces of $\mathbb{F}_2^n$, denoted by 
    ${n\brack \ell}$, and two vertices $L$ and $L'$ are adjacent if 
    ${\sf dim}(L\cap L') = \ell-1$. 
\end{enumerate}
Both of the graphs above, namely the Johnson graph with large intersection sizes, as well as the Grassmann graph, are not 
small set expanders. However very importantly, the class
of small sets in the Grassmann graph with bad expansion has a succinct and intuitive characterization, 
and the proof of the $2$-to-$1$ Games Theorem heavily relies on this characterization. 

% To get some intuition for this characterization, notice 
% that the Grassmann graph has a collection
% of induced subgraphs isomorphic to lower order 
% Grassmann graphs, each forming a set with bad expansion. These can be described by the two basic sets,
% \[
% Z_x = \left\{L\in {n\brack \ell}~\Big|~x\in L\right\},
% \qquad
% Z_W = \left\{L\in {n\brack \ell}~\Big|~L\subseteq W\right\}
% \]
% where $x\in\mathbb{F}_2^n\setminus\{0\}$ is a point and $W\subseteq \mathbb{F}_2^n$ is a hyperplane. The sets $Z_x$ and $Z_W$ are often
% thought of as local sets, since membership in them is determined 
% by a single vector (either belonging to $L$ or to its orthogonal complement).
% It is easy to observe that $Z_x$ and $Z_W$ both have edge expansion at most $1/2$, and taking intersection of any $r$ of them yields a set that has expansion at most $1-2^{-r}$. 
% In essence, the aforementioned characterization asserts that a 
% small set $S$ in the Grassmann graph that has expansion at most $1-\eps$ is correlated with an intersection of some $r$-basic sets, where $r = \log(1/\eps)$. 

Though the term is not formally defined, we refer to graphs such 
as the Grassmann graph above as globally hypercontractive graphs. 
By that, we mean that there is a collection of 
``obviously-non-expanding local sets'', such that any small set that 
doesn't expand well must have a large intersection with one of the sets from 
the collection (see Section~\ref{sec:abstraction} for a semi-formal definition). Aside from the Grassmann graph, this 
class of graphs includes Johnson graphs with small intersection sizes~\cite{KMMS},
certain Cayley graphs over 
the symmetric group~\cite{filmus2020hypercontractivity},
$p$-biased cubes for $p=o(1)$, other product domains~\cite{keevash2021global} as well as high dimensional expanders~\cite{gur2022hypercontractivity,bafna2022hypercontractivity}.


A natural question to consider then refers to the 
complexity of UG over graphs that are
globally hypercontractive, which is the main question 
of study in this paper. The motivation for this 
is two-fold:
\begin{enumerate}
    \item First, in light of 
the reduction of~\cite{KhotMS17,DBLP:conf/stoc/DinurKKMS18a,DinurKKMS18,KhotMS18} and the above discussion, it is an interesting 
question whether there may be a reduction showing the NP-hardness of UG that only uses globally hypercontractive graphs. Gaining a better understanding for these questions will shed further
light on the source of hardness in this reduction, as well as explain the distinct roles of the ``smooth parallel repetition'' and ``composition with the Grassmann encoding'' steps above.

\item Second, the UGC is related to another well-known computational problem known as the Small-set Expansion (SSE) problem~\cite{RaghavendraS10,RaghavendraST12}. In the SSE problem one is given a graph $G = (V,E)$ such that either (a) for all $S\subseteq V$ of size at most $\gamma|V|$ it holds that $\Phi(S)\geq 1-\eps$, or (b) there is $S\subseteq V$ of size at most $\gamma|V|$ such that $\Phi(S)\leq \delta$. The
Small-set Expansion hypothesis (SSEH) asserts that for all $\eps,\delta>0$
this problem is NP-hard for sufficiently small $\gamma$, and the 
papers~\cite{RaghavendraS10,RaghavendraST12} show that this 
hypothesis implies the UGC.\footnote{To be more precise, it implies
a stronger version of the UGC, which states that UG 
remains hard on graphs that have an expansion profile
as the noisy hypercube.} 

Despite recent progress towards UGC, no new results for the SSEH were proved. In some sense, it seems that to prove hardness results for SSE, one would need 
to only use gadgets that are small-set expanders or perhaps globally hypercontractive. Otherwise, in the final graph one may create small non-expanding sets that are unintended (namely, that don't correspond to any solution to the 
initial problem we reduced from). This motivates the study of the role of globally hypercontractive graphs in reductions related to Unique-Games.
\end{enumerate}

The question has been addressed recently in~\cite{BBKSS} (with 
follow-up by~\cite{DBLP:conf/soda/BafnaHKL22}) for the class of \emph{Affine 
Unique-Games}, which we define and discuss in the next section.

\subsection{Unique-Games over Globally Hypercontractive Graphs}
\begin{definition}
An instance of $\text{Affine-UG}$ is an instance of Unique-Games 
in which the alphabet is the ring of integers modulo $q$, $\mathbb{Z}_q$, and all of the 
constraint maps $\phi_e$ are affine shifts, that is, $\phi_e$ 
of the form $\phi_e(\sigma) = \sigma+b_e$ for some $b_e\in\mathbb{Z}_q$.
\end{definition}
An equivalent but slightly different way to view 
the Affine-UG problem is as a system
of linear equations $(X,E)$ over $\mathbb{Z}_q$.  
Each equation in $E$ is of the form $x_i - x_j = b$ where $x_i,x_j\in X$ are variables 
and $b\in\mathbb{Z}_q$ is some constant. Despite looking very restrictive, 
it is known~\cite{KhotKMO07} that the UGC is true if and 
only if it holds for the class of Affine UG and furthermore this class captures many interesting optimization problems such as Max-Cut and graph coloring, thus we shall focus our attention on Affine UG henceforth. 
\footnote{We remark that the 
reduction of~\cite{KhotKMO07} does not preserve the topology of the graph. 
We are therefore not able to translate our results directly to the 
class of general UG, and believe this is an interesting direction 
for further study.}

In~\cite{BBKSS}, the authors investigate the complexity of Affine UG 
on the Johnson graph, albeit with small intersection sizes, the regime in which a characterization theorem for non-expanding sets holds. 
Namely, they consider
the ``$\alpha$-noisy'' Johnson graph $J(n,\ell,t)$ in the case that $t = \alpha \ell$, for $\alpha\in (0,1)$ bounded away from $0$ and thought of as a fixed constant independent of $\ell$.
The result of~\cite{BBKSS} 
asserts that for small enough $\eps>0$, there is a polynomial time algorithm
that given an instance of UG over $J(n,\ell,t)$ with value at least 
$1-\eps$, finds an assignment that satisfies at least $C(\ell,\alpha,\eps)>0$
of the constraints. The most important feature of this result is the lack
of dependency on the size of the graph $n$ and the alphabet size, however 
it has two significant downsides:
% Here, we consider
% the Johnson graph $J(n,\ell,t)$ in the case that $t = \alpha \ell$ for $\alpha\in (0,1)$ thought of as bounded away from $1$. In~\cite{KMMS}, it is shown
% that small sets in $J(n,\ell,t)$ that have edge expansion at most $1-\eps$ 
% are dense inside the induced subgraph on a set of the form
% \[
% H_R = \left\{A\in \binom{n}{\ell}~\Big|~A\supseteq R\right\},
% \]
% where $|R|\leq O\left(\frac{\log(1/\eps)}{\alpha}\right)$. 
% Indeed, the sets $H_R$ are the analogs of the basic sets for the setting of the Johnson graph, and we will also refer to them by this name. 
% Similarly 
% to before, the induced subgraph on any basic set $H_R$ is isomorphic to a lower order
% Johnson graph $J(n-r,\ell-r,t-r)$ where $r = |R|$. The result of~\cite{BBKSS} 
% asserts that for small enough $\eps>0$, there is a polynomial time algorithm
% that given an instance of Unique-Games over $J(n,\ell,t)$ with value at least 
% $1-\eps$, finds an assignment that satisfies at least $C(\ell,\alpha,\eps)>0$
% of the constraints. The most important feature of this result is the lack
% of dependency on the size of the graph $n$ and the alphabet size, however 
% there are two significant downsides to it:
\begin{enumerate}
    \item The completeness guarantee needed for the algorithm to work is close to $1$ ($\eps$ close to $0$). Thus, within the realm of their result, it is still possible that UG over the noisy-Johnson and Grassmann graphs is hard for completeness smaller than $1/2$ (wherein 
    general UG is now known to be hard) but suddenly becomes easy when the completeness exceeds $1/2$. In fact, UG with completeness $> 1/2$ is easy for the Grassmann graph via the UG algorithm for expanders itself, thus their result does not imply anything non-trivial about the complexity of UG on the Grassmann graph, despite its being globally hypercontractive. Therefore, it is not
    clear what their algorithm implies about the hard instances of~\cite{DinurKKMS18} which have completeness less than $1/2$: (1) are they hard due to the smooth parallel repetition step, or the composition with the Grassmann graph? (2) Is the smoothness of the outer PCP necessary -- instead of the Johnson graph with large intersection sizes could one have instead used the noisy-Johnson graph?    
    \item Their soundness guarantee depends on $\ell$, the uniformity of the sets in the Johnson graph, which in most PCP constructions grows with the alphabet size of the instance. Indeed, such graphs are often used to encode a global function $f\colon [n]\to\Gamma$ (often from an error correcting code) where $|\Gamma| = O(1)$. 
    In that case, the Johnson graph corresponds to the restrictions encoding
    $F[A] = f|_{A}$, in which case each vertex $A$ of the Johnson graph is 
    assigned a label from $\Gamma^{\ell}$ corresponding to the restriction of 
    $f$ to $A$. Thus, in the hard instances of UG obtained via PCP constructions (in particular the one above) the parameter $\ell$ 
    grows with the alphabet size of the UG instance, hence we would like the performance of our algorithm to not depend on it. 
    
    We remark though, that it is okay for the runtime of our algorithm to depend on it (just like an algorithm 
    running in time $n^{k}$ that solves UG 
    would refute the conjecture, but the trivial random algorithm that satisfies at least a $1/k$ fraction of the constraints does not).
\end{enumerate}

\subsection{Our Results}
Our main result asserts that there is a polynomial time algorithm for solving 
Affine UG over globally hypercontractive graphs that bypasses these 
two barriers thus addressing the motivations presented in Section~\ref{sec:ug}. As the term globally hypercontractive graph is not formally defined, 
below are some concrete instances of graphs on which this applies. In Section~\ref{sec:abstraction} we give a semi-formal definition of globally hypercontractive graphs and also show how our algorithm and analysis can be abstracted to solve UG on such graphs, as long as one is provided with an SoS certificate of global hypercontractivity.

We first consider the noisy-Johnson graph, 
for which we have the following two results.
The first result addresses the case that 
the completeness of the instance is close to 
$1$, in which case our algorithm matches the 
guarantee of the algorithm of~\cite{BBKSS} 
for certifiable small-set expanders, and in particular the $\alpha$-noisy hypercube graph:
\begin{theorem}\label{thm:main_johnson_close_to_1}
There is $\eps_0>0$
such that for all $\alpha\in (0,1)$ 
the following holds for all $0<\eps\leq \eps_0$.
There exists an algorithm whose running time is $n^{\poly(\ell,|\Sigma|,1/\eps)}$ which, on input
$\Psi$ which is an affine UG instance over $J(n,\ell,\alpha\ell)$ promised to be at least $(1-\eps)$-satisfiable, finds an assignment that satisfies 
at least $2^{-O\left(\frac{\sqrt{\eps}}{\alpha}\right)}$-fraction of the constraints in $\Psi$.
\end{theorem}


The second result addresses the case of UG instances with arbitrarily small (but bounded away from $0$) completeness, in which case our algorithm satisfies a constant fraction of the constraints:
\begin{theorem}\label{thm:main_johnson}
For all $\alpha\in (0,1)$ and $c>0$, there is $\delta>0$ such that the following holds. 
There exists an algorithm whose running time is $n^{D}$ with $D = \ell^{\poly(|\Sigma|\ell^{1/c})}$\footnote{We note that we have not optimized for $D$ and the $\exp(\ell)$-dependence arises due to the degree of the SoS proofs in Section~\ref{sec:johnson_2}. We used a blackbox statement to convert some of the proofs therein to SoS proofs, and we conjecture that one can in fact improve the SoS degree to $O(\ell)$ when done carefully.} which on input
$\Psi$, an affine UG instance over $J(n,\ell,\alpha \ell)$ promised to be at least $c$-satisfiable, finds an assignment that satisfies 
at least $\delta$-fraction of the constraints in $\Psi$.
\end{theorem}

We note that the soundness guarantee in the theorems above does not depend on $\ell$ when $\alpha = \Omega(1)$, but degrades as $\alpha$ approaches $0$. In Section~\ref{sec:interpretation} below we discuss why one cannot get a polytime algorithm that obtains an $\Omega(1)$-valued solution when $\alpha = o(1)$ (e.g. $1/\sqrt{\ell}$) therefore showing that this dependence on $\alpha$ is necessary.

We can get similar results given any of the globally hypercontractive graphs mentioned earlier. Below
we give a corollary for the Grassmann graph.
We show that there is a polynomial time algorithm solving affine UG over the Grassmann graph, even on 
instances with small completeness:
\begin{theorem}\label{thm:main_grass}
For all $c>0$ there exists $\delta>0$ such that the following holds. There exists an algorithm whose running time is $n^{D}$ with $D = \ell^{\poly(|\Sigma|\ell^{1/c})}$ which
on input $\Psi$, an affine UG instance over 
${\sf Grass}(n,\ell)$ promised to be at least $c$-satisfiable, finds an assignment that satisfies 
at least $\delta$-fraction of the constraints in $\Psi$.
\end{theorem}

As discussed before, since the spectral gap of the Grassmann graph is $1/2$, UG algorithms over expanders already imply Theorem~\ref{thm:main_grass} for $c \gg 1/2$. Thus, the main contribution of Theorem~\ref{thm:main_grass} is the algorithm on Grassmann graphs that works for arbitrarily small completeness.

Below we state our result for random walks on high dimensional expanders (HDX), a large class of graphs that generalize the Johnson graphs but do not necessarily possess its strong symmetries. These  include graphs stemming from cut-offs of \cite{lubotzky2005explicit}'s construction of Ramanujan complexes, or  \cite{kaufman2018construction}'s construction of coset complex expanders. These graphs exhibit the nice high-dimensional expansion properties (e.g. global hypercontractivity) of the Johnson graphs yet are substantially different in other aspects, such as being of bounded degree.

\begin{theorem}\label{thm:main_hdx}
For all $\alpha \in (0,1)$ and $c>0$, there exists $\delta > 0$ such that the following holds. Let $X$ be any $d$-dimensional two-sided $\gamma$-local-spectral expander with $\gamma \ll o_\ell(1)$ and $d>\ell$. There exists an algorithm whose running time is $n^{D}$ with $D = \ell^{\poly(|\Sigma|\ell^{1/c})}$ which on input $\Psi$, an affine UG instance over the canonical walk $M$ on $X(\ell)$ of depth $\alpha$, promised to be at least $c$-satisfiable, finds an assignment that satisfies 
at least $\delta$-fraction of the constraints in $\Psi$.
\end{theorem}

Since we have not defined any of the HDX terminology, let us note that this is indeed a generalization of Theorem~\ref{thm:main_johnson}. The Johnson graph corresponds to the complete complex $X$ (which is the simplest instantiation of a two-sided local spectral expander), and the $\alpha$-noisy Johnson graph $J(n,\ell,\alpha\ell)$ corresponds to a ``canonical'' random-walk on $X(\ell)$ that goes down $\alpha\ell$-levels and comes back up randomly to $X(\ell)$ while ensuring that it changes exactly $\alpha\ell$ elements in a vertex. In fact, in the above theorem we can allow $M$ to be any complete random walk on $X(\ell)$ and our soundness guarantee will only depend on $c$ and certain parameters of $M$ that are inherently independent of $\ell$ \footnote{Concretely it depends on the stripped threshold rank of $M$ above a certain threshold as defined in~\cite{DBLP:conf/soda/BafnaHKL22}. For example, when $M$ is the canonical random walk with depth $\alpha$ on $X(\ell)$, and the completeness is $c=1-\eps$, this quantity is $r(M) = O(\sqrt{\eps}/\alpha)$ and our soundness guarantee is $\exp(-r)$, matching that of Theorem~\ref{thm:main_johnson_close_to_1}.}. 


For the sake 
of simplicity of presentation, in the first part of the paper we focus
on the Johnson graph. In Section~\ref{sec:other_graphs} we provide an abstraction of our techniques and discuss the algorithm for general globally hypercontractive graphs, in particular for the Grassmann graph and random-walks over HDXs.

%We remark that our techniques can be adjusted so as to get analogous algorithms for other globally hypercontractive graphs such as the ones mentioned above. 

\subsubsection*{An Interpretation of Our Results}\label{sec:interpretation}
Our results suggest that the hardness in 
the instances of UG obtained via the reduction  of~\cite{KhotMS17,DBLP:conf/stoc/DinurKKMS18a,DinurKKMS18,KhotMS18} does not come from
the Grassmann graph (which is globally 
hypercontractive), but rather from the smooth parallel repetition step. Recall that this step uses a Johnson graph with a large intersection
parameter ($J(n,\ell,\alpha\ell)$ with $\alpha \approx 0$), that is not globally-hypercontractive. Therefore combining the knowledge from the reduction and our algorithm we get that the $\alpha$-noisy-Johnson graphs are hard for UG when $\alpha = o(1)$ and become easy when $\alpha$ is bounded away from $0$, thus also explaining why our soundness guarantee must necessarily decay with $\alpha$ (under $P \neq NP$). Indeed, we would be able to make such an assertion provided that our
results held for general UG (as opposed to only affine UG) or if the reduction above produced  instances of Affine UG. Though
we believe an algorithm for general UG should exist along the lines of our algorithm, we 
do not know how to prove so and leave this is 
an interesting direction to investigate. 

Albeit, ignoring the subtlety between general and affine UG, this means that any future progress on UGC will have to use graphs that are \emph{not} 
globally hypercontractive, possibly again 
via the technique of \emph{smooth} parallel repetition. 

\subsection{Techniques: New rounding scheme for higher degree SoS}
Our algorithms are obtained via a novel rounding scheme and analysis for the standard higher degree Sum-of-Squares SDP relaxation for Unique Games. Raghavendra's~\cite{Raghavendra08} groundbreaking result showing the optimality of the basic SDP for all CSPs under the UGC, led to efforts to refute the UGC using higher degree SoS relaxations~\cite{Lasserre00,Parrilo00}. The study of SoS algorithms has since produced numerous algorithmic advances across many fronts: high-dimensional robust statistics~\cite{BarakKS14,ma2016polynomial, BakshiDHKKK20, BakshiDJKKV22}, quantum computation~\cite{BKS17} and algorithms for semi-random models~\cite{BuhaiKS22}, to name a few. Most of these works use the sum-of-squares method for \emph{average-case} problems though and unfortunately there remains a dearth of techniques for analysing higher degree SoS relaxations for \emph{worst-case} optimization problems. The handful of techniques known for worst-case rounding are the \emph{global correlation rounding} technique from~\cite{BarakRS11,RT12} and its generalization via reweightings in~\cite{BKS17}.

In this context, very recently~\cite{BBKSS} proposed a new technique for rounding relaxations of UG that have ``low-entropy'' measured via a function called the shift-partition size. Given two fixed assignments for the instance, their shift-partition size is roughly defined as the fraction of variables on which these assignments agree (upto symmetry).
Taking the equivalent view of the SDP solution as a distribution $\cD$ over non-integral solutions, called a pseudodistribution, the expected shift-partition size of two random assignments drawn from $\cD$ is then roughly equal to an average of local collision probabilities under $\cD$ and thus a proxy for the \emph{entropy} of $\cD$. Their analysis proceeds by showing: (1) when the expected shift-partition size (equivalently collision probability) is large, one can round to a high-valued solution, and moreover (2) when the graph is a certifiable small-set expander, the pseudodistribution always has large shift-partition size! They were not able to extend this idea to get high-valued solutions for the broader class of globally hypercontractive graphs though, since in this case the pseudodistribution might be supported over multiple assignments and therefore does not have high collision probability. It turns out though that even in this harder case, the pseudodistribution $\cD \times \cD$ has large expected shift-partition size after \emph{conditioning} on an event $E$. But they could not exploit this property since after conditioning the shift-partition could be large for trivial reasons\footnote{In the worst case, the event $E$ could collapse the product distribution over two random assignments to set the second random assignment to be always equal to the first one. In this case a pair of assignments drawn from $\cD \times \cD ~|~ E$ being equal does not say anything about the collision probability of $\cD$.} and therefore is no longer a good proxy for the collision probability/entropy of the distribution $\cD$.

Our main technical contribution is to strengthen and broadly extend this framework of rounding low-entropy pseudodistributions. We show that after a suitable preprocessing step on the pseudodistribution, one can in fact condition on any event $E$ (with not too small probability) while preserving most of the desired local independence properties of the distribution. Thus, even after conditioning on $E$, the expected shift-partition size of $\cD\times\cD ~|~ E$ being large signifies that the pseudodistribution $\cD$ has high collision probability. One can then use a simple rounding procedure to obtain a high-valued UG solution. Conditioning pseudodistributions is one of the few ways we know of harnessing the power of higher-degree pseudodistributions, hence we believe that the idea of gaining structural control over the distribution after conditioning may be applicable in the analysis of other SoS algorithms too.

At a high level, the algorithm of~\cite{BBKSS} can be viewed as showing one how to round a pseudodistribution when it is essentially supported over \emph{one} assignment (upto symmetry of the solution space). In general though, our problem might not be as structured and the resulting pseudodistribution for it might not be low-entropy in this restricted sense\footnote{For example, it could be supported over $O(1)$ solutions, in which case the collision probability would be high only after conditioning on an appropriate event.}. We expand this notion of ``low-entropy'' thus allowing a broader set of pseudodistributions, intuitively those with a ``few good solutions'', to fall into this class and therefore become amenable to rounding. We expect that with this strengthening, this framework  should be broadly applicable for algorithm design for other optimization problems. Below is a detailed overview of our techniques, starting out with the framework of~\cite{BBKSS}.



\begin{comment}
Currently there are only a handful of techniques available to analyze sum-of-squares relaxations for worst-case optimization problems. For low-degree SoS programs (namely degree $2$), 
one typically uses the vector-valued solution
in order to generate Gaussian samples with a
given covariance matrix, which can then be 
translated into an integral solution using 
rounding~\cite{GW, sparsest-cut}. In a higher degree sum-of-squares 
program, one is also able to condition on
the integral value of a few variables in
the program, and after that still have a 
valid sum-of-squares solution. Typically, 
after conditioning one no longer has a good
structural control over the vector-valued 
solution, and the only thing that can really be done at that point is a Gaussian rounding scheme as above. 


\mitali{We refine and strengthen the framework of rounding via the shift-partition technique.}

The lack of structural control after conditioning is precisely the bottleneck
in the approach of~\cite{BBKSS}. Therein,
the authors consider an even $E$ which has
noticeable (but not too large) probability
which is defined by a low-degree polynomial.
Ideally, they would like to be able to condition on the event $E$ as well as retain various nice structural properties their original solution has. If this was possible, the guarantees of the algorithm of~\cite{BBKSS} would be much stronger and would not deteriorate with the graph size. 
Alas, as they are 
unable to do so, they have to give up
conditioning on the event $E$, and as a 
result the probability of $E$ gets absorbed
into the soundness guarantee of their algorithm.

As our main conceptual contribution, 
we show that after a suitable pre-processing step, one can in fact condition on 
an event $E$ as above while keeping 
various good global properties of 
the solution. Below is a detailed overview of our techniques.
% as done in~\cite{BBKSS}. In this section, we 
% give a high level overview of the approach laid down therein, and then explain how we build on top of it to prove our results.

\end{comment}

\subsubsection{The Approach of~\cite{BBKSS}: Rounding analysis via the Shift Partition}
Fix an Affine Unique-Games instance $\Psi = (G=(V,E), \mathbb{F}_q, \Phi)$. 
In the SoS relaxation of the Unique-Games problem we have a collection 
of variables $X_{v,\sigma}$, one for pair of vertex $v\in V$ and label 
to it $\sigma\in \Sigma$. The output of the program is a pseudoexpectation 
operator $\widetilde{\E}$, which assigns to each monomial involving 
at most $d$ of the variables a real-number, under which:
\begin{enumerate}
    \item The value is high:
    \[
    \widetilde{\E}\left[\sum\limits_{(u,v)\in E}\sum\limits_{\sigma\in \Sigma} X_{v,\sigma} X_{u,\phi_{u,v}(\sigma)}\right]\geq c\cdot|E|.
    \]
    \item $\widetilde{\E}$ is a linear, positive semi-definite operator (when viewed as a matrix over $\mathbb{R}^{M\times M}$ where $M$ is the set of monomials of degree at most $d/2$) satisfying various 
    Booleanity constraints on $X_{u,\sigma}$.
    \item Scaling: $\widetilde{\E}[1] = 1$.
\end{enumerate}
Morally, the pseudoexpectation 
$\widetilde{\E}$ should be thought of 
in the following way: 
there is an unknown distribution $\mathcal{D}$ 
over assignments $A_1,\ldots,A_m$ that each have value 
at least $c$. For the 
assignment $A_i$ we think of Boolean valued assignment to 
the variables $X_{u,\sigma}$ that assigns to a variable $1$
if and only if $A_i(u) = \sigma$, and associate with it 
the expectation operator $\E_i$ which maps monomials to 
Boolean values in the natural way according to $A_i$. The operator 
$\widetilde{\E}$ then is the average of the operators $\E_i$ 
according to $i\sim \mathcal{D}$. 
\footnote{Formally speaking,
when given $\widetilde{\E}$ we are not guaranteed that there 
exists an actual distribution $\mathcal{D}$ over good assignments 
as above, however this intuition will be good enough 
for the sake of this informal presentation.}

\paragraph{Shift-partition:}
Given $\widetilde{\E}$, one can construct a different pseudoexpectation operator that allows access 
to moments of two assignments 
$X = A_i, X' = A_j$ where $i,j\sim \mathcal{D}$ are 
chosen independently. 
In expectation, we get that at least $c^2$ fraction of the edges
get satisfies by both $X$ and $X'$; the algorithm attempts 
to satisfy these edges. Towards this end, given two fixed assignments $X$ and $X'$ we define the shift-partition of the vertices of $V$: $V = \cup_{s\in\mathbb{F}_q} F_s$ where for each $s\in\mathbb{F}_q$ we define
\[
F_s(X,X') = \left\{v\in V~|~X(v) - X'(v) = s\right\}.
\]
The shift-partition size is then defined as: 
\[\pE_{X,X' \sim \cD}\left[\sum_{s \in \Sigma} \left(\frac{|F_s(X,X')|}{|V(G)|}\right)^2\right].\]  
After rearranging, we get that when $X$ and $X'$ are independent, this expression is an average of some local collision probabilities (precisely $\E_{u,v}[CP(X_u - X_v)]$), and hence the shift-partition size being large in expectation turns out to be useful for rounding.

On the other hand, observe that if an edge $(u,v)\in E$ is satisfied by both $X$ and 
$X'$, then $X(u) - X(v) = X'(u) - X'(v)$ and rearranging we 
conclude that $u$ and $v$ are in the same part $F_s$ 
of the shift partition. We therefore conclude that in expectation over $X,X' \sim \cD$ at least $c^2$ fraction of the edges of $G$ stay
inside the same part of the shift partition, implying that the expansion of the shift-partition is small. 

\paragraph{Small-set expanders.} If the graph 
$G$ is a small-set expander, then the above implies that at
least one of the sets $F_s$/the shift-partition size is large and the following 
rounding procedure works in such cases:
\begin{enumerate}
    \item Sample a vertex $v\in V$ and choose $A(v) = \sigma$ according to the distribution $p(\sigma) = \widetilde{\E}[X_{v,\sigma}]$.
    \item For any $u\in V$, sample $A(u)$ according to the distribution $p(a)=\frac{\widetilde{\E}[X_{u,a}X_{v,\sigma}]}
    {\widetilde{\E}[X_{v,\sigma}]}$.
\end{enumerate}
To get an understanding to why this rounding scheme works, think of $X$ as fixed and $X'$ as random. Thus, the fact that part $s$ of
the shift partition is large implies 
that $X' = X+s$ 
on a constant fraction of the vertices.
Therefore, once we sampled the assignment to $v$ in the first part 
of the algorithm, the value of $s$ is determined. In the second step we are sampling the assignment to other 
nodes conditioned on the value of $v$. However, there is one
value for $u$ which is much more likely than others -- namely 
$X(u)+s$, and so we can expect that $X'(u) = X(u) + s$ for a constant fraction of the vertices $u$. In particular, for any edge $(u,w)$ inside $F_s$ that is satisfied by $X$, we will have that the assignments sampled for $u$ and $w$ are $X(u)+s$ and $X(w) + s$ respectively with constant probability, in which case we manage to satisfy $(u,w)$. To analyse this rounding strategy formally,~\cite{BBKSS} crucially use the independence of $X$ and $X'$. 
%since in this case we get that the expected shift-partition size is equal to a weighted collision probability of $\cD$.

In essence, the above asserts that the shift-partition being large implies that the solution space of $X$ must have high collision probability, which can then be used for rounding. By that, we mean that our distribution essentially consists of only one assignment (upto shift-symmetry) and its perturbations.

\paragraph{Non small-set expanders.} Consider a graph which is 
not a small set expander, say that $G$ is the Johnson graph $J(n,\ell,t=\ell/2)$. In that case the above reasoning no longer 
works as $F_s$ may indeed be all small sets. However, as explained 
earlier, using global hypercontractivity we can infer that one of the sets $F_s$ must posses a 
certain structure -- it must have large density inside one of the canonical non-expanding sets.
In the case of the Johnson graph specifically, these canonical sets take the following form:
\[
H_R = \left\{A\in\binom{n}{\ell}~\Big|~A\supseteq R\right\},
\]
for $R \subset [n]$ and $|R| = r = O(1)$. In fact, global hypercontractivity gives the following stronger structural property: the set $H = \bigcup_{R\in\mathcal{R}} H_R$ where $\mathcal{R}$ 
consists of all $R$'s inside which some part $F_s$ is dense, 
has a constant measure. Doing simple accounting, it follows that $|\mathcal{R}|\geq \Omega(n^r/\ell^r)$
and as there are at most $\binom{n}{r}$ different canonical sets it follows that $|\mathcal{R}|$ contains an
$\Omega(1/\ell^r)$ fraction of these sets.

For each choice of $X$ and $X'$ though 
we may have a different collection of dense subcubes $\mathcal{R}$. But since $\mathcal{R}$ contains an $\Omega(1/\ell^r)$ fraction of all the subcubes, we get that there must be at least one subcube $H_R$ that is dense with probability $\Omega(1/\ell^r)$ over $X,X' \sim \cD$. Let $H_R$ be such a subcube and $E_R(X,X')$ be the event that $H_R$ is dense. Ideally, at this point one would like to condition on $E_R$ so that one of the parts inside the shift partition $F_s$ becomes large inside $H_R$, and then hope that as was the case for small-set expanders, we can satisfy many of the edges inside $F_s\cap H_R$. 

%letting $E_R$ be the event that $F_s$ is dense inside $H_R$ for some $s$, we get that on average over $X,X' \sim \cD$ at least
%\[
%\frac{|\mathcal{R}|}{\binom{n}{r}}
%\geq \Omega\left(\frac{1}{\ell^r}\right)
%\]
%of the events $E_R$ hold. 


%and then recurse (since still, the fraction of edges inside $F_s\cap H_R$ is relatively small, but this is because this is a small set of vertices). 
Unfortunately, this hope does not materialize -- after conditioning on $E_R$ even though the shift-partition is large, the rounding strategy above may break. Indeed, 
for the rounding procedure we wanted the values of $X(u)$ and $X'(u)$ for $X,X'\sim\mathcal{D}$ to be independent 
for every vertex $u$. However, after conditioning the joint distribution $\cD \times \cD ~|~E_R$ over $(X,X')$ might have correlations between $X$ and $X'$. In particular this distribution could even be supported on pairs $(X,X')$ that are always equal to each other, in which case the shift-partition is large because of trivial reasons and therefore its large size doesn't imply anything about the collision probability/entropy of $\cD$. 

Hence in~\cite{BBKSS} the authors don't manage to do this conditioning, and instead settle 
for satisfying an $\Omega\left(\frac{1}{\ell^{2r}}\right)$-fraction of the constraints on $H_R$. After that they iterate this algorithm many times to satisfy an $\Omega\left(\frac{1}{\ell^{2r}}\right)$-fraction of the constraints of the whole graph. 
%They show that one can do without conditioning and still satisfy at least  fraction of the edges. 


\subsubsection{Our Approach: Conditioning on the Event $E$ via
(Eliminating) Global Correlations}
Our main contribution to the 
above framework is to show 
that by adding an additional preprocessing step, we can ensure that even after conditioning on the event $E_R$ 
above, the assignments $X$ and 
$X'$ will remain highly 
independent. In particular, 
the fact that some part in
the shift partition becomes
large must happen -- just like
in the case of small-set expanders -- due to the fact
that our distribution has high collision probability.

As the event $E = E_R(X,X')$ has 
probability at least $\Omega\left(\frac{1}{\ell^r}\right)$, 
 if we are sufficiently high up in the SoS hierarchy 
($\Theta(\ell^r)$ levels will do, for an overall running time of 
$n^{\Theta(\ell^r)}$), we do have access to the conditional pseudoexpectation
\[
\pE[Y~|~E] = \frac{\pE[Y 1_E]}{\pE[1_E]}.
\]
This means that we can sample labels of vertices 
conditioned on the event $E$. To make this useful though, we must change the 
rounding procedure. 
To get some intuition consider 
the extreme case in which after conditioning on $E(X,X')$ 
there are huge correlations between $X$
and $X'$ that remain in our distribution. 

Namely, suppose that after conditioning on $E$ it holds that $X(u) = X'(u)$ for almost all vertices $u$. 
In that case, if we sampled 
$X,X'$ from $\cD \times \cD$ (not conditioned on $E$), we would get that with probability 
at least $\Pr[E]\geq \Omega\left(\frac{1}{\ell^r}\right)$
the event $E$ holds, in which case $X$ and $X'$ agree on almost all vertices. This means that if $\cD$ was an actual distribution the assignments have a large global correlation: fix $X' = X_0$ for $X_0$ that satisfies $\Pr_\cD[E(X,X_0) = 1] \geq \Pr_{\cD \times \cD}[E]$. Once $E$ holds, 
we have that $X(u) - X(v) = X_0(u) - X_0(v)$ for almost all pairs of vertices, hence the values of the assignment $X$ to the vertices $u$ and $v$ is correlated across $\cD$.
Therefore, a natural idea is to avoid this issue by transforming $\cD$ to another distribution lacking global correlations, in the sense that the assignments to a typical pair of vertices $u$ and $v$ are almost independent. 


For this purpose we use an idea from~\cite{RT12}, which adapted to our setting says 
that for any $\tau>0$ there is $d = d(\tau,|\Sigma|)$ such that conditioning $\pE$ on the values of $d$ randomly chosen vertices ensures that the global correlation is 
at most $\tau$. That is, the values of $X(u)$ and $X(v)$
for two typical vertices $u$ and $v$ are at most $\tau$-correlated, and the same holds for $X'$. 
In Lemma~\ref{lem:correlation} we then show that if we start with such a pseudodistribution that lacks global correlations, then one can condition on the event $E$ and retain near independence between the assignments $X$ and $X'$, at least on most vertices. To be more precise, we show that for $Y_{u,v} = (X(u), X(v))$ and $Y_{u,v}' = (X'(u), X'(v))$, the statistical distance between $Y_{u,v}, Y'_{u,v}~|~E$ 
and $Y_{u,v}, Y'_{u,v}$ is small for almost all pairs of vertices $u,v$.\footnote{To make our rounding succeed we need to use a more complicated version of $Y_{u,v}$ (see Definition~\ref{defn:local-dist}). Proving Lemma~\ref{lem:correlation} then turns out to be  technically challenging since these $Y_{u,v}$'s are not indicator variables, only approximately so and additionally we need to make sure that our proof works for pseudodistributions instead of just for actual distributions.}

Using this idea we are able to get an $\Omega(1)$-valued solution on some basic set $H_R$. To summarize, we first preprocess the pseudodistribution to eliminate global correlations. We can then find an event $E(X,X')$, corresponding to the 
fact that some part $F_s$ in the shift partition has becomes dense in some basic set $H_R$. Furthermore, conditioning on $E$ most pairs $(X(u),X(v)), (X'(u),X'(v))$ remain almost-independent. Then running a simple rounding procedure on $H_R$ (as in~\cite{BBKSS}), we are able to satisfy a good fraction of the edges inside $H_R$. $H_R$ might be a $o(1)$ fraction of the graph though, therefore like~\cite{BBKSS} we repeat this procedure multiple times to get an $\Omega(1)$-valued solution for the whole graph. This gives an efficient algorithm for affine UG over the Johnson graphs as in Theorem~\ref{thm:main_johnson_close_to_1}.

To prove Theorem~\ref{thm:main_johnson}
(namely, the regime where $c$ is not close to $1$) more work is needed. Indeed, in the case that $c$ is 
close to 
$1$ we are able to conclude that essentially all edges 
stay within some part $F_s$ of the shift partition. Thus, 
as long as our sets $H_R$ cover a constant fraction of 
the edges that stay within some $F_s$, they are automatically guaranteed to cover a constant fraction of the edges that 
are satisfied by both $X$ and $X'$, and these are the edges our rounding procedure manages to satisfy. If $c$ is just 
bounded away from $0$ we can no longer make such an argument, 
and it is no longer even clear that the sets 
$H_R$ cover some edges that we have a hope of satisfying.


\subsubsection{Getting Small Completeness: Capturing all of the Non-expanding Edges}
To design our algorithm for the case when the completeness $c$ is just 
guaranteed to be bounded away from $0$ we must first argue that in the shift partition, we are able to capture almost all 
of the edges that stay within a part $F_s$  
using the basic sets $H_R$ (so as to ensure we are including 
the edges that $X$ and $X'$ both satisfy).

Towards this end we require a more
refined corollary of global hypercontractivity, asserting 
that if we have a small set of vertices $F$ in the Johnson 
graph that has edge expansion at most $1-\eta$, then we 
can find a collection $\mathcal{R}$ of basic sets such that:
\begin{enumerate}
    \item \textbf{Bounded and dense:} 
    each $R\in \mathcal{R}$ has size $|R| = O(1)$ 
    and $F$ is dense inside each $H_R$. That is, 
    $\delta(F\cap H_R)\geq \Omega_{\eta}(\delta(H_R))$ for 
    each $R\in\mathcal{R}$.
    \item \textbf{Maximally dense:} 
    For all $R\in\mathcal{R}$ and all $R'\subsetneq R$, 
    $F$ is not very dense in $H_{R'}$.
    \item \textbf{Capture almost all non-expanding edges:} 
    Almost all the edges that stay inside $F$ also stay inside 
    $H_R$ for some $R\in\mathcal{R}$.
\end{enumerate}
Indeed, we show that a global hypercontractive inequality such as the one 
in~\cite{KMMS} can be used to prove such a result (in a black-box manner). 


Using this result, we are able to argue that that the edges that 
stay inside the subcubes $H_R$ for $R\in\mathcal{R}$  
cover most of the edges that stay within the same part in the shift
partition. There are several subtleties here that one has to deal with, for example, ``regularity issues'' such as, how many different $R$'s cover 
a given edge. The goal of the second item above is to handle such concerns, and it roughly says that no vertex nor 
edge gets over-counted by a lot. After that, we are able to 
condition on an event $E$, where as before $E$ indicates that some part 
$F_s$ becomes dense inside some basic set $H_R$, so that the resulting distribution has a large shift-partition inside $H_R$.
At this point, we are (morally) 
back to the problem of rounding the SoS solution 
on a set with a large shift-partition, except that now our solution 
has value $c'>0$ (as opposed to close to $1$). We remark that again, we use the ``elimination of global correlations'' 
idea presented earlier to retain near independence after conditioning. With more care, we use a similar analysis to the one presented for completeness close to $1$ to finish the proof when $c$ is arbitrarily small.



\subsection{Open Problems}
We end this introductory section by stating a few open
directions that are of interest for future research.
The first problem asks whether our results continue to 
hold for non-affine unique games:
\begin{problem}
For globally hypercontractive graphs $G$ 
such as the Johnson graph 
(with small intersection size) and the Grassmann graph, is there
a polynomial time algorithm that given a UG instance 
$\Psi$ over $G$ with ${\sf val}(\Psi)\geq 1-\eps$ (where $\eps>0$ is thought of as small), finds 
an assignment satisfying at least $\delta$ fraction of the
constraints in $\Psi$? How about the case that ${\sf val}(\Psi)\geq c$, where $c$ is bounded away from $1$?
\end{problem}

The second problem asks whether there are other combinatorial optimization problems for which our techniques may yield improved algorithms. Informally, we show how to round pseudodistributions with low-entropy, including for instance those that are supported over $O(1)$ assignments (and their perturbations). We do so by proving that after the
elimination of global correlations one can retain local independence properties even after conditioning. We believe that this technique should be useful outside the context of UG -- given any problem for which one can prove (in SoS) that there are only a ``few good solutions'', one can apply our rounding technique to obtain one such solution.

%At the heart of our algorithm lies the fact that using the elimination of global correlations one can retain pairwise-independent type properties even after conditioning, and we believe this observation should be useful outside the context of UG.
\begin{problem}
Can one use the low-entropy rounding framework to get improved run-time for other combinatorial optimization problems, such as coloring $3$-colorable graphs using as few colors as possible?
%Can one use the elimination of global correlation + rounding techniques to get improved run-time over other combinatorial optimization problems, such as coloring $3$-colorable graphs using as few colors as possible?
\end{problem}

Third, it would be interesting to see whether our techniques 
can help in designing improved sub-exponential time algorithms
for combinatorial optimization problems such as Max-Cut over general graphs (ideally) or over special classes of graphs.
\begin{problem}
Can one use the low-entropy rounding framework to get improved approximation algorithm for Max-Cut that runs in time $2^{n^{\eps}}$? 
% \mitali{as written it would violate UGC right? i think we want $2^{n^\eps}$-time algorithms for $(1-\eps,1-O(\eps))$-Max cut. Probably this is called near-exponential?}
\end{problem}






